\documentclass[11pt]{article}
\usepackage[a4paper,margin=1in]{geometry}
\usepackage{mathptmx}
\usepackage[hidelinks]{hyperref}
\pagestyle{empty}
\setlength{\parskip}{4pt}
\setlength{\parindent}{0pt}

\begin{document}

\begin{flushright}
October 22, 2026
\end{flushright}

Editor-in-Chief\\
IEEE Transactions on Intelligent Transportation Systems

\vspace{6pt}
Dear Editor-in-Chief,

I submit the manuscript ``Multi-Hover-Point Relay Scheduling for UAV
Emergency Delivery under a Fixed Transport Plan'' for consideration as a
regular paper in IEEE Transactions on Intelligent Transportation Systems.

Mountain floods damage roads and communication facilities at the same time, so
emergency delivery and communication support have to be organized together
under blocked roads, tight energy budgets and incomplete information. The
manuscript builds a ``transport scheduling--communication support--resource
allocation'' chain on one consistent physical model: energy-aware batching
(18 trips, 59.24\,kWh, 0.69\% below FFD/BFD on the test instance),
heterogeneous-fleet routing with battery and time windows (20 trips,
65.25\,kWh, 80 boxes and all 31 hard deadlines met), task-driven relay
scheduling, and partition-based resource accounting. Every headline number is
recomputed from the frozen artifacts by an independent read-only verifier
(relay-required count within 0.3\%, relay energy within 0.1\%).

The central contribution is the multi-hover-point relay scheduling mode
(MHPRS): one relay UAV serves several hover points per take-off/return cycle
instead of returning to base after each service window. On the test instance,
two relay airframes cover all 11,713 relay-required samples on the 1-s grid
with four sorties at 5.43\,kWh (joint energy 70.68\,kWh), whereas the tested
single-hover heuristics reach only 51.0\%, 62.7\% and 82.1\% with one, two and
three airframes; the single-hover feasible set is a subset of the proposed one.
The mode is formulated under a fixed transport plan and coupled to DEM-based
cruise altitudes, energy and airframe/battery turnover constraints, which is
where it differs from relay-placement or trajectory studies that optimize a
single sortie. The manuscript explicitly does not claim a joint optimum: the four stages run
in a dependency order, and the paper states this positioning and the scope of
its numerical claims.

I also wish to be transparent about the instance. It is a reconstructed
relief scenario: two of the seven service windows re-derive hover points from
the DEM grid rather than read them from the original log, so the agreement is
reported as a reconstruction check, not as a flight test. The conventions used
in the reconstruction---an aggregate 37\,s climb allowance for A/B departures,
a raised corridor on S007 legs, and a 6\,m clearance with bilinear
interpolation---are declared in the limitations section. The 100\% coverage
figure is tied to the stated verification convention; under a stricter 6\,m
check one reconstructed hover point misses its link threshold by 0.96\,dB, and
a local 256\,m adjustment restores sampled coverage of the affected window with
582\,s of arrival slack.

The manuscript is original, has not been published in a peer-reviewed venue
and is not under consideration elsewhere. I declare no competing interests;
the work received no specific funding. Code, instance data, the benchmark suite and the read-only
verification pipeline are available at
\url{https://github.com/vscode-aiyaya/uav-flood-transport}.

\vspace{10pt}
Sincerely,

\vspace{16pt}
Xin Yang\\
School of Mechanical and Power Engineering, Zhengzhou University\\
Zhengzhou, Henan 450001, China\\
\href{mailto:xinyang.zzu@outlook.com}{xinyang.zzu@outlook.com}

\end{document}
